\section{Casimir element}

\begin{definition}
    Let $L$ be a finite-dimensional Lie algebra over $F$, and let
    $B(\cdot,\cdot)$ be a nondegenerate invariant symmetric bilinear form on $L$.
    Choose a basis $x_1,\dots,x_n$ of $L$, and let $x^1,\dots,x^n$ be the
    $B$-dual basis, so that
    \[
        B(x_i,x^j)=\delta_{ij}.
    \]
    The element
    \[
        \Omega_B=\sum_{i=1}^n x_i x^i \in U(L)
    \]
    is called the Casimir element associated with $B$.
\end{definition}

\begin{remark}
    The definition is independent of the choice of basis. Indeed, $\Omega$ is the element corresponding to the identity map of $L$ under the natural identification
    \[
        L \otimes L \cong \operatorname{End}_F(L),\qquad
        x \otimes y \mapsto \bigl(z \mapsto (y,z)x\bigr).
    \]
\end{remark}

\begin{proposition}
    Let $L$ be a finite-dimensional Lie algebra equipped with a nondegenerate invariant symmetric bilinear form $B$, and let
    \[
        \Omega_B=\sum_{i=1}^n x_i x^i
    \]
    be the associated Casimir element. If $\rho:L\to\mathfrak{gl}(V)$ is a representation of $L$, define
    \[
        C_V:=\sum_{i=1}^n \rho(x_i)\rho(x^i)\in \operatorname{End}_F(V).
    \]
    Then
    \[
        \rho(a)C_V=C_V\rho(a)
        \qquad\text{for all } a\in L.
    \]
\end{proposition}

\begin{proof}
    For any $a\in L$, we compute
    \[
        [\rho(a),C_V]
        =
        \sum_{i=1}^n \bigl([\rho(a),\rho(x_i)]\rho(x^i)+\rho(x_i)[\rho(a),\rho(x^i)]\bigr).
    \]
    Since $\rho$ is a representation, this becomes
    \[
        [\rho(a),C_V]
        =
        \sum_{i=1}^n \bigl(\rho([a,x_i])\rho(x^i)+\rho(x_i)\rho([a,x^i])\bigr).
    \]

    Write
    \[
        [a,x_i]=\sum_{j=1}^n c_{ji}x_j.
    \]
    We claim that
    \[
        [a,x^i]=-\sum_{j=1}^n c_{ij}x^j.
    \]
    Indeed, for every $k$,
    \[
        B([a,x^i],x_k)=-B(x^i,[a,x_k])
    \]
    by invariance and symmetry of $B$. Since
    \[
        [a,x_k]=\sum_{j=1}^n c_{jk}x_j,
    \]
    we get
    \[
        B([a,x^i],x_k)
        =
        -\sum_{j=1}^n c_{jk}B(x^i,x_j)
        =
        -c_{ik}.
    \]
    If
    \[
        [a,x^i]=\sum_{j=1}^n d_{ji}x^j,
    \]
    then
    \[
        d_{ki}=B([a,x^i],x_k)=-c_{ik},
    \]
    because $B(x^j,x_k)=\delta_{jk}$. This proves the claim.

    Therefore
    \[
        \sum_{i=1}^n \rho([a,x_i])\rho(x^i)
        =
        \sum_{i,j=1}^n c_{ji}\rho(x_j)\rho(x^i),
    \]
    while
    \[
        \sum_{i=1}^n \rho(x_i)\rho([a,x^i])
        =
        -\sum_{i,j=1}^n c_{ij}\rho(x_i)\rho(x^j).
    \]
    After relabeling indices in the first sum, these two sums cancel. Hence
    \[
        [\rho(a),C_V]=0.
    \]
    Equivalently,
    \[
        \rho(a)C_V=C_V\rho(a)
        \qquad\text{for all } a\in L.
    \]
\end{proof}

